\section{结论}

本文提出了 \textbf{EdgeSched}，这是一个面向内存受限消费级设备本地 AI 运行时的轻量级资源调度框架。通过正视专用显存约束，EdgeSched 使一个常驻服务与一个突发式 GPU 密集任务能够在同一台设备上共存，而无需依赖昂贵的硬件升级。

本文的主要贡献包括：
\begin{enumerate}
\item \textbf{不对称资源迁移}：将用于快速抢占的 ``Urgent Mode''（GPU→CPU）与用于低中断恢复的 ``Atomic Hot-Swap''（CPU→GPU）结合起来，在显存安全与服务连续性之间取得平衡。

\item \textbf{轻量级持久化}：基于 mmap 的软件定义状态背板，以很小的单次操作开销提供崩溃一致性（在我们的测量中，注册/更新开销低于 0.01 ms），并有助于避免系统失败后积累僵尸任务。

\item \textbf{可预测的 QoS}：通过优先级队列、反压机制和取消语义，使系统在测试的饱和场景下保持响应性与稳定性。
\end{enumerate}

在消费级硬件（RTX 5070 Laptop）上的实验结果支持了该方法在所评估混合工作负载场景中的有效性。当更简单的基线面临 OOM 风险或反复重载代价时，EdgeSched 能在突发任务执行期间保持常驻服务可达，并在崩溃后约 44 ms 内将任务状态恢复到语义上可用的状态。虽然其原始吞吐略低于简单信号量方案（在我们的 mock 微基准中约低 16\%），但在可靠性与连续性上的收益使其成为稳健本地 AI 应用的可行方案。

\subsection{未来工作}

未来工作将沿以下方向展开：

\begin{enumerate}
\item \textbf{异构多 GPU 支持}：扩展 EdgeSched 以支持多 GPU 配置，从而实现更细粒度的资源分配和负载均衡。

\item \textbf{内核级调度}：集成更细粒度的内核级调度，以实现更低时延的抢占并进一步缩短服务中断时间。

\item \textbf{自适应策略}：根据历史负载模式和用户行为，动态调整迁移策略与队列容量，实现更智能的资源管理。

\item \textbf{更多模态支持}：扩展系统以支持文本转语音（TTS）、视觉语言模型（VL）和其他模态任务，构建完整的多模态 AI 助手生态。
\end{enumerate}

EdgeSched 为本地消费级设备上的多模态 AI 应用提供了一种务实的调度方案。我们希望这项工作能够推动关于紧资源约束下运行时管理的进一步研究。